\section{Retrospective lifecycle-trajectory diagnosis}
\label{sec:v5-trajectories}
This analysis reuses all nine accepted fixed-H physical runs and the twelve
unique O/C/C30/E predictions from the frozen V3 evaluation: three windows,
three physical repeats per window, and four model comparisons per run.
The resulting 36 pairings are not new experiments or independent prediction
samples. We neither refit a model nor replace the original M1 endpoint.
The artifact binds the consumed records to their original acceptance hashes
and rechecks all 44,100 decisions across the 21 unique traces.

\subsection{Definitions and event alignment}
Time is relative to each window start, on $[0,2100)$ seconds.
Physical occupied capacity counts each generation from launch until verified
release, clipped to that interval; simulated capacity includes every non-off
state, including stopping. Let $n_m(t)$ and $n_r(t)$ be predicted and observed
occupied GPU counts. For model $m$ and repeat $r$, we distinguish
\begin{align}
A_{m,r} &= \left|\int_0^{2100}(n_m(t)-n_r(t))\,dt\right|,\\
T_{m,r} &= \int_0^{2100}|n_m(t)-n_r(t)|\,dt.
\end{align}
Both have units GPU\,$\cdot$\,s. The first is the existing total-occupation
absolute error; the second measures trajectory discrepancy without signed
temporal cancellation. We integrate the piecewise-constant records exactly,
without substituting one-second samples for continuous occupation.
Table~\ref{tab:v5-groups} reports equal-repeat means within each window;
Table~\ref{tab:v5-pairs} preserves every pairing.

Targets are compared at all 2,100 one-second decision ticks. The first
state-count difference compares the decision-time counts of active, starting,
draining/stopping, and off devices; it is not the first difference of all
hidden system state. Queue and running observations precede each decision.
Routes are compared over all arrival IDs as local, synthetic API, or
undispatched. For a request whose final routes differ, the recorded first
dispatch evidence is a retrospective locator, not a time when both eventual
routes were already known. Route differences need not follow target differences.

\subsection{An endpoint discrepancy at the cooldown boundary}
The frozen C/C30 configurations differ only in the startup endpoint.
Their first common starts and subsequent divergences are shown in
Table~\ref{tab:v5-endpoints}. C30 becomes ready after 30 s, but targets stay
equal until the 60 s cooldown expires. At that tick C remains starting;
its queue-plus-running pressure exceeds the strict up-threshold of eight,
and its target rises from one to two. C30 is ready with lower pressure and
keeps target one. Both controllers have identical target/cooldown state
before this decision; swapping the observed inputs swaps this one-step action.
This diagnoses the frozen simulator contrast, not a new physical intervention
or a replay of an alternative future trajectory. At the corresponding tick,
each of the nine physical repeats has one active device, no starting device,
and target one; that is an observational anchor.
H does not read the predictor's startup parameter, and the guard's separately
frozen safety-delay parameters are not intervened on here.

\begin{table*}[t]
\centering
\small
\caption{First C/C30 startup contrast in all three windows. Times are seconds
from window start; pressure is queued plus locally running requests. The first
target difference occurs at common start plus 60 s in each window.}
\label{tab:v5-endpoints}
\par\vspace{4pt}
\begin{tabular*}{0.92\linewidth}{@{\extracolsep{\fill}}lrrrrr@{}}
\toprule
Window & Common start & C30 ready & C ready & First target difference & Pressure C/C30\\
\midrule
% Generated from accepted V5 retrospective analysis; do not edit values.
Burst & 195 & 225 & 264.219 & 255 & 13/1\\
Recovery & 126 & 156 & 195.219 & 186 & 14/2\\
Steady & 171 & 201 & 240.219 & 231 & 21/4\\
\bottomrule

\end{tabular*}
\end{table*}

\subsection{Counterexamples and temporal cancellation}
The six burst/recovery physical repeats first differ from all four models
in target at $t=3$ s, before any window restart. At that tick the models
observe no local running request while the physical controllers still do;
completion observation has crossed a decision tick. These differences cannot
be attributed to window-restart timing. Conversely, steady repeats 1 and 3 have exactly
the same 2,100 targets as C30 despite 154 and 152 state-count difference
ticks, 48 and 45 route differences, and nonzero occupied-resource errors.
Readiness error therefore does not necessarily change a target.

The steady E total MAE is 66.34 GPU\,$\cdot$\,s, close to C30's 67.18,
but its mean trajectory error is 1,390.50 versus 114.32 GPU\,$\cdot$\,s.
The same E trace disagrees on 66.73\% of target ticks on average.
Accurate total occupation can thus hide substantial temporal mismatch;
this diagnostic does not retroactively select a different model or alter M1.
All C30 pairings retain nonzero occupation error. The largest C30 target
residual is burst repeat 2 (637/2100 ticks); its longest mismatch spans
$[1020,1144)$ s. This example is selected by an explicitly post-hoc maximum
residual rule after the full audit, alongside all no-target-change cases.
The artifact exports every window's first common-start segment for all models
and repeats, as well as the complete startup and unmatched-pair counts.

\begin{table*}[t]
\centering
\small
\caption{All window/model means over three physical repeats. Total MAE is
mean $A_{m,r}$; trajectory error is mean $T_{m,r}$, both in GPU\,$\cdot$\,s.
Target disagreement uses 2,100 ticks per pair; route disagreement uses all
arrival IDs in that window, including undispatched requests.}
\label{tab:v5-groups}
\par\vspace{4pt}
\begin{tabular*}{0.92\linewidth}{@{\extracolsep{\fill}}llrrrr@{}}
\toprule
Window & Model & Total MAE & Trajectory error & Target diff. (\%) & Route diff. (\%)\\
\midrule
% Generated from accepted V5 retrospective analysis; do not edit values.
Burst & O & 234.48 & 457.24 & 15.17 & 27.95\\
Burst & C & 847.60 & 1369.86 & 49.63 & 27.61\\
Burst & C30 & 83.08 & 268.96 & 10.33 & 15.26\\
Burst & E & 70.55 & 237.76 & 10.33 & 15.71\\
Recovery & O & 244.08 & 421.31 & 11.71 & 27.75\\
Recovery & C & 864.20 & 1302.19 & 43.92 & 35.24\\
Recovery & C30 & 78.30 & 138.31 & 3.24 & 2.20\\
Recovery & E & 66.95 & 99.02 & 3.24 & 4.19\\
Steady & O & 1151.97 & 1517.83 & 59.40 & 42.84\\
Steady & C & 964.14 & 1621.50 & 73.68 & 38.87\\
Steady & C30 & 67.18 & 114.32 & 1.03 & 4.99\\
Steady & E & 66.34 & 1390.50 & 66.73 & 49.45\\
\bottomrule

\end{tabular*}
\end{table*}

\begin{table*}[t]
\centering
\small
\caption{All 36 model/physical-repeat pairings. Errors are GPU\,$\cdot$\,s;
target differences are counts out of 2,100 ticks. A dash means targets never
differ, not a missing record. Route counts retain the full arrival denominator.
The twelve unique predictions are reused across the nine physical runs.}
\label{tab:v5-pairs}
\par\vspace{4pt}
\begin{tabular*}{0.92\linewidth}{@{\extracolsep{\fill}}llrrrrrr@{}}
\toprule
Window & Model & Rep. & Total error $A$ & Trajectory $T$ & Target diff. & First diff. (s) & Route diff./arrivals\\
\midrule
% Generated from accepted V5 retrospective analysis; do not edit values.
Burst & O & 1 & 300.54 & 304.39 & 147 & 3 & 65/297\\
Burst & O & 2 & 103.59 & 764.07 & 662 & 3 & 117/297\\
Burst & O & 3 & 299.30 & 303.26 & 147 & 3 & 67/297\\
Burst & C & 1 & 913.66 & 1347.35 & 1056 & 3 & 76/297\\
Burst & C & 2 & 716.71 & 1416.14 & 1015 & 3 & 90/297\\
Burst & C & 3 & 912.42 & 1346.08 & 1056 & 3 & 80/297\\
Burst & C30 & 1 & 53.54 & 53.54 & 7 & 3 & 2/297\\
Burst & C30 & 2 & 143.41 & 701.03 & 637 & 3 & 128/297\\
Burst & C30 & 3 & 52.30 & 52.30 & 7 & 3 & 6/297\\
Burst & E & 1 & 4.48 & 6.53 & 7 & 3 & 6/297\\
Burst & E & 2 & 201.43 & 698.53 & 637 & 3 & 122/297\\
Burst & E & 3 & 5.72 & 8.24 & 7 & 3 & 12/297\\
Recovery & O & 1 & 305.21 & 335.21 & 187 & 3 & 122/454\\
Recovery & O & 2 & 122.16 & 593.86 & 364 & 3 & 133/454\\
Recovery & O & 3 & 304.85 & 334.85 & 187 & 3 & 123/454\\
Recovery & C & 1 & 925.33 & 1233.56 & 895 & 3 & 156/454\\
Recovery & C & 2 & 742.28 & 1440.00 & 977 & 3 & 167/454\\
Recovery & C & 3 & 924.97 & 1233.01 & 895 & 3 & 157/454\\
Recovery & C30 & 1 & 52.21 & 52.21 & 9 & 3 & 10/454\\
Recovery & C30 & 2 & 130.84 & 310.86 & 186 & 3 & 11/454\\
Recovery & C30 & 3 & 51.85 & 51.85 & 9 & 3 & 9/454\\
Recovery & E & 1 & 5.81 & 8.75 & 9 & 3 & 21/454\\
Recovery & E & 2 & 188.86 & 279.36 & 186 & 3 & 16/454\\
Recovery & E & 3 & 6.17 & 8.95 & 9 & 3 & 20/454\\
Steady & O & 1 & 1193.59 & 1538.83 & 1255 & 99 & 686/1564\\
Steady & O & 2 & 1068.73 & 1475.71 & 1232 & 99 & 635/1564\\
Steady & O & 3 & 1193.58 & 1538.95 & 1255 & 99 & 689/1564\\
Steady & C & 1 & 1005.76 & 1623.09 & 1548 & 231 & 596/1564\\
Steady & C & 2 & 880.91 & 1618.33 & 1546 & 231 & 635/1564\\
Steady & C & 3 & 1005.76 & 1623.10 & 1548 & 231 & 593/1564\\
Steady & C30 & 1 & 76.70 & 76.70 & 0 & -- & 48/1564\\
Steady & C30 & 2 & 48.15 & 189.57 & 65 & 969 & 141/1564\\
Steady & C30 & 3 & 76.70 & 76.70 & 0 & -- & 45/1564\\
Steady & E & 1 & 74.18 & 1388.10 & 1400 & 232 & 784/1564\\
Steady & E & 2 & 50.68 & 1395.18 & 1404 & 232 & 755/1564\\
Steady & E & 3 & 74.17 & 1388.23 & 1400 & 232 & 781/1564\\
\bottomrule

\end{tabular*}
\end{table*}

These observations support checking lifecycle states, actions, and resource
trajectories alongside matched-service fit. Three controlled windows and nine
repeats do not establish a universal startup/cooldown law, randomized physical
causality, or a general guard guarantee. Software-observed readiness is not a
GPU-kernel timestamp, and the synthetic API retains its original quality and
billing limitations.
